\documentclass[dvipdfmx]{jsarticle}
\usepackage{tikz,ascmac,amsmath,amssymb,amsthm,mathtools,fancybox,bm,physics,arydshln,mathcomp,wrapfig}
\usetikzlibrary{calc}
\usetikzlibrary{intersections}

\begin{document}
\begin{center}
 \ 
 
\vspace{15zw}
  {\HUGE レポート課題}
  
\vspace{3zw}
{\huge HRID: 336229}

\bigskip

{\huge 三森誠真}

\end{center}
\newpage

1. 因数分解せよ(係数は整数のみに限る)

(1) $x^2 - 2x -2{,}400$

2次方程式$x^2 - 2x -2{,}400 = 0$ を解くと, \ 

$$
\begin{aligned}
  x &= 1 \pm \sqrt{1 + 2{,}400} \\
  &= 1 \pm \sqrt{2{,}401} \\
  &= 1 \pm 49 \\
  &= -48, \ 50
\end{aligned}
$$

であるから, \ \underline{$x^2 - 2x -2{,}400 = (x + 48)(x - 50).$}

\bigskip

(2) $12x^2 -41x + 24$

$$
\begin{aligned}
  12x^2 -41x + 24 
  &= (3 \cdot 4)x^2 +\{3\cdot(-3) + 4\cdot(-8)\}x + (-8)\cdot(-3) \\
  &= \underline{(3x - 8)(4x - 3)}.
\end{aligned}
$$

\bigskip

(3) $4x^2 - y^2 + 6yz - 9z^2$

$$
\begin{aligned}
  4x^2 - y^2 + 6yz - 9z^2
  &= 4x^2 - (y^2 - 6yz + 9z^2) \\
  &= 4x^2 - (y-3z)^2 \\
  &= \underline{(2x + y - 3z)(2x - y + 3z)}.
\end{aligned}
$$

\bigskip

(4) $x^3 + 14x^2 + 56x + 64$

$$
\begin{aligned}
  x^3 + 14x^2 + 56x + 64
  &= (x+4)(x^2 + 10x + 16) \ \ (\because 因数定理) \\
  &= \underline{(x+2)(x+4)(x+8)}.
\end{aligned}
$$

\bigskip

(5) $x^6 - y^6$

$$
\begin{aligned}
  x^6 - y^6
  &= (x^3 + y^3)(x^3 - y^3) \\
  &= (x+y)(x^2 -xy - y^2)(x-y)(x^2 + xy + y^2) \\
  &= \underline{(x+y)(x-y)(x^2 + xy + y^2)(x^2 -xy - y^2)}
\end{aligned}
$$

\bigskip

(6) $(x + y + z)^3 -x^3 - y^3 - z^3$

$$
\begin{aligned}
  x^3 + y^3 + z^3
  &= \{(x+y)^3 -3xy(x+y)\} + z^3 \\
  &= \{(x+y)^3 + z^3\} -3xy(x+y) \\
  &= (x+y+z)^3 - 3z(x+y)(x+y+z) - 3xy(x+y)
\end{aligned}
$$

であるから, \ 
$$
\begin{aligned}
  (x + y + z)^3 -x^3 - y^3 - z^3
  &= 3z(x+y)(x+y+z) + 3xy(x+y) \\
  &= 3(x+y)(z^2 + xy + yz + zx) \\
  &= \underline{3(x+y)(y+z)(z+x)}.
\end{aligned}
$$

(7) $x^2 + 4xy + 4y^2 + 5x + 10y + 6$

$$
\begin{aligned}
  x^2 + 4xy + 4y^2 + 5x + 10y + 6
  &= (x+2y)^2 + 5(x+2y) + 6 \\
  &= \underline{(x+2y+2)(x+2y+3)}.
\end{aligned}
$$

\bigskip

(8) $x^2y + y^2z - y^3 - x^2z$

$$
\begin{aligned}
  x^2y + y^2z - y^3 - x^2z
  &= x^2(y-z) - y^2(y-z) \\
  &= \underline{(x+y)(x-y)(y-z)}.
\end{aligned}
$$

\bigskip

(9) $x^3 - y^3 -3xy -1$

$$
\begin{aligned}
  x^3 - y^3 -3xy -1
  &= x^3 + (-y)^3 + (-1)^3 -3\cdot x \cdot (-y) \cdot (-1)\\
  &= (x -y - 1)(x^2 + y^2 + (-1)^2 +xy - y + x) \\
  &= \underline{(x- y - 1)(x^2 + y^2 + xy + x - y + 1)}.
\end{aligned}
$$

\bigskip

(10) $x^5 - x^4 + x^3 - x^2 + x -1$

$$
\begin{aligned}
  x^5 - x^4 + x^3 - x^2 + x -1
  &= (x-1)(x^4 + x^2 + 1) \ \ (\because 因数定理) \\
  &= (x-1)\{(x^4 + 2x^2 + 1) - x^2 \} \\
  &= (x-1)\{(x^2 + 1)^2 - x^2 \} \\
  &= \underline{(x-1)(x^2 + x + 1)(x^2 - x + 1)}.
\end{aligned}
$$

\bigskip

(11) $(x+1)(x+2)(x+3)(x+6) - 3x^2$
$$
\begin{aligned}
  (x+1)(x+2)(x+3)(x+6) - 3x^2
  &= (x+1)(x+6)(x+2)(x+3) - 3x^2 \\
  &= (x^2 + 7x + 6)(x^2 + 5x + 6) -3x^2 \\
  &= (x^2 + 6 + 7x)(x^2 + 6 + 5x) - 3x^2 \\
  &= (x^2 + 6)^2 + 12x(x^2 + 6) +32x^2 \\
  &= (x^2 + 6 + 4x)(x^2 + 6 + 8x) \\
  &= \underline{(x^2 + 4x + 6)(x^2 + 8x + 6)}.
\end{aligned}
$$

\bigskip

(12) $(x+y)(y+z)(z+x) +xyz$

$$
\begin{aligned}
  (x+y)(y+z)(z+x) +xyz
  &= \{x^2 + (y+z)x + yz\}(y+z) + xyz \\
  &= (y+z)x^2 + \{(y+z)^2 + yz\}x + yz(y+z) \ \ (xについての降べき順) \\
  &= (x+y+z)\{(y+z)x  + yz\} \\
  &= \underline{(x+y+z)(xy + yz + zx)}.
\end{aligned}
$$


\newpage

2.

(1) $a^x - a^{-x} = 5^{0.5}$ のとき $a^x + a^{-x}, \ a^{3x} + a^{-3x}$ の値を求めよ

\bigskip

$a^{2x} + a^{-2x} = (a^x - a^{-x})^2 + 2a^x a^{-x} = 5 + 2 = 7$ であるから, \ 

$(a^x + a^{-x})^2 = a^{2x} + a^{-2x} + 2 = 9$.

$a^x + a^{-x} > 0$ であるから, \ \underline{$a^x + a^{-x} = 3$}.

$a^{3x} + a^{-3x} = (a^x + a^{-x})(a^{2x} + a^{-2x} - 1) = 3(7-1) = \underline{18}$.

\bigskip
\bigskip

(2) $2^x + 2^y = 160, \ 2^{x+y} = 4{,}096, \ x<y$ を満たす$x, y$ を求めよ

\bigskip

$2^x, \ 2^y$ は二次方程式$X^2 - 160X + 4{,}096 = 0$ の2解であるから, \ まずこの二次方程式を解く.

$4{,}096 = 2^{12}$ であることに着目すると, \ 
$X^2 - 160X + 4{,}096 = (X - 2^5)(X-2^7)$.

よって, \ $X^2 - 160X + 4{,}096 = 0$ を解くと, \ $X = 2^5, \ 2^7$.

このとき, \ $x < y$ より, \ $2^x = 2^5, \ 2^y = 2^7$ が得られる.

したがって, \ \underline{$x=5, \ y=7$}.

\bigskip
\bigskip


\begin{wrapfigure}{r}{0pt} %図形を右側に
\begin{tikzpicture}[scale=1]
  % 軸
  \draw[->] (-1,0) -- (3,0) node[right] {$x$};
  \draw[->] (0,-1) -- (0,5) node[above] {$y$};


  % 直線 y = -2x + 4 の破線（x<0 と x>2 の部分）
  \draw[dashed,domain=-0.5:0,smooth,variable=\x] plot ({\x},{-2*\x+4});
  \draw[dashed,domain=2:2.5,smooth,variable=\x] plot ({\x},{-2*\x+4});

  % 実線（0 <= x <= 2）
  \draw[thick,domain=0:2,smooth,variable=\x] plot ({\x},{-2*\x+4});

  %関数名
  \node[scale=0.8, right] at (1.5,2) {$y=-2x+4$};


  % 交点のマークとラベル（グラフ上に表示）
  \fill (2,0) circle (1.5pt) node[below right, xshift=2pt, yshift=-2pt] {2};
  \fill (0,4) circle (1.5pt) node[above left,  xshift=-2pt, yshift=2pt]  {4};
  \fill (0,0) circle (1.5pt) node[below left, xshift=2pt, yshift=2pt] {O};
\end{tikzpicture}
\end{wrapfigure}




(3) $2x + y = 4, \ x \ge 0, \ y \ge 0$ のとき $a^{xy} (a > 0)$ の最大値を求めよ

\bigskip

$2x + y = 4, \ x \ge 0, \ y \ge 0$ という条件は, \ $y = -2x + 4, \ 0 \le x \le 2$ と言い換えることができ(右図参照), \ 

\begin{wrapfigure}[2]{r}{0pt} %図形を右側に
\begin{tikzpicture}[scale=1]
  % 軸
  \draw[->] (-1,0) -- (3,0) node[right] {$x$};
  \draw[->] (0,-3) -- (0,3) node[above] {$y$};

  % 直線 y = -2x + 4 の破線（x<0 と x>2 の部分）
  \draw[dashed,domain=-0.5:0,smooth,variable=\x] plot ({\x},{-2*pow(\x,2) + 4*\x});
  \draw[dashed,domain=2:2.5,smooth,variable=\x] plot ({\x},{-2*\x^2+4*\x});

   % 実線（0 <= x <= 2）
  \draw[thick,domain=0:2,smooth,variable=\x] plot ({\x},{-2*\x^2+4*\x});

  %関数名
  \node[scale=0.8, right] at (2,1) {$y=-2(x-1)^2+2$};

  % 交点のマークとラベル（グラフ上に表示）
  \fill (1,0) circle (1.5pt) node[below , yshift=-1pt] {1};
  \fill (2,0) circle (1.5pt) node[below right, xshift=2pt, yshift=-2pt] {2};
  \fill (0,2) circle (1.5pt) node[above left,  xshift=-2pt, yshift=2pt]  {2};
  \fill (0,0) circle (1.5pt) node[below left, xshift=2pt, yshift=2pt] {O};

  %点線
  \draw[dashed] (1,0)--(1,2);
  \draw[dashed] (0,2)--(1,2);
\end{tikzpicture}
\end{wrapfigure}


$a^{xy} = a^{-2x^2 + 4x} \ (a > 0)$ において,

(i) $0 < a < 1$ のとき,

$a^{xy}$ は単調減少するため, \ 関数$f(x) = -2x^2 + 4x$ の$0 \le x \le 2$における最小値を求めればよい.


具体的に$x=0, \ 2$ のとき$f(x)$ は最小となり, \ 最小値は$0$ である.

よって, \ このとき$a^{xy}$ は最大となり, \ 最大値は$a^0 = 1$ となる.

\bigskip

(ii) $a = 1$ のとき.

$a^{xy} = 1$ であるから最大値も$1$.

\bigskip

(iii) $a > 1$ のとき,
 
$a^{xy}$ は単調増加するため, \ 関数$f(x) = -2x^2 + 4x$ の$0 \le x \le 2$における最大値を求めればよい.

具体的に$x=1$ のとき$f(x)$ は最大となり, \ 最大値は$2$ である.

よって, \ このとき$a^{xy}$ も最大となり, \ 最大値は$a^2$ となる.

\bigskip

(i) ～ (iii) より,

$$\underline{
 \begin{dcases}
    0 < a \le 1 のとき, & 最大値は1 \\
    a > 1 のとき,       & 最大値はa^2 
  \end{dcases}
  }
$$

\newpage

(4) 第 $10$ 項が $210$, \ 第$20$項が$-10$ である等差数列の初項から第$n$ 項までの和を$S_n$ とするとき, \ $S_n$ の最大値を求めよ

\bigskip

問題の数列を$\{a_n\}$ とし, \ 公差を$d$ とする. このとき, \ 条件より
$210 + 10d = -10$. よって$d = -22$.

初項$a_1$ については, \ $a_1 + 9d = a_{10}$ より, \ 
$a_1 = 210 - 9d = 210 + 198 = 408$.

$$
\begin{aligned}
  よって, \hspace{5zw} S_n &= \dfrac{(a_1 + a_n)n}{2} \\
  &= \dfrac{n}{2}\{2a_1 + (n-1)d \} \\
  &= -11n^2 + 419n \\
  &= -11\left(n - \dfrac{419}{22} \right)^2 + \dfrac{419^2 \cdot 11}{22^2}
\end{aligned}
$$

実数全体を定義域にもつ関数$f(x) = -11\left(x - \dfrac{419}{22} \right)^2 + \dfrac{419^2 \cdot 11}{22^2}$ を考えたときに 

$y = f(x)$ のグラフは$x = \dfrac{419}{22}$ に関して対称であり, \ 
$n$ は自然数, \ かつ $\dfrac{419}{22} \fallingdotseq 19.05$ であることから, \ 

$S_n$ は$n=19$ において最大となり, \ 最大値は$S_{19} = -11\left(19 - \dfrac{419}{22} \right)^2 + \dfrac{419^2 \cdot 11}{22^2} = \underline{3{,}990}$.

\bigskip
\bigskip

(5) 初項$1$, \ 公比$(1+r)^{-1}$ の等比数列の第$n$項から第$2n$項までの和を求めよ

\bigskip

条件より, \ 第$n$項は$\{(1+r)^{-1}\}^{n-1} = (1+r)^{-(n-1)}$ であるため, \ 第$n$項から第$2n$項までの和は, \ 

$$
\begin{aligned}
  \dfrac{(1+r)^{-(n-1)}\{1 - \{(1+r)^{-1}\}^{n+1} \}}{1 - (1+r)^{-1}}
  &= \dfrac{(1+r)^{n+1} - 1}{(1+r)^{2n} - (1+r)^{2n-1}} \ \ (分母分子を(1+r)^{2n}倍 ) \\
  &= \underline{\dfrac{(1+r)^{n+1} - 1}{r(1+r)^{2n-1}}}.
\end{aligned}
$$

\bigskip
\bigskip

3.

(1) mutimaro という文字列の文字すべてを用いてできる順列の個数を求めよ

\bigskip

重複する文字mを区別して並び方を考えると$8!$ 通り. 

そして, \ 文字mの並び方$2$通りを考慮すると, \ $\dfrac{8!}{2} = \underline{20{,}160}$

\bigskip
\bigskip

(2) $7$色の玉をつないでブレスレットを作る方法は何通りあるか

\bigskip

$7$個の数珠順列となるため, \ $\dfrac{(7-1)!}{2} = \underline{360 \ (通り)}$

\bigskip
\bigskip

(3) 正$n$角形の頂点を結んでできる三角形の個数はいくつか

\bigskip

$n$個の頂点のうち$3$つを選ぶ組み合わせを考えればよいため, \ 
${}_n \mathrm{C}_3 = \underline{\dfrac{n(n-1)(n-2)}{6}}$.

\newpage

(4) コインを$10$回投げて表が奇数回出る場合は何通りあるか

\bigskip

表が出る回数は$1, \ 3, \ 5, \ 7, \ 9$ 回のどれかであるため, \ それぞれの回数に対して, \ 

表が出るコインを選ぶ組合せを考えると, \ 

${}_{10} \mathrm{C}_1 + {}_{10} \mathrm{C}_3 + {}_{10} \mathrm{C}_5 + {}_{10} \mathrm{C}_7 + {}_{10} \mathrm{C}_9 = \underline{512 \ (通り)}$

\bigskip
\bigskip

(5) $x, \ y, \ z$ は $x + y + z = 20$ を満たす$0$以上の整数とする. このような$(x, \ y, \ z)$ の組は何組あるか

\bigskip

まず, \ $x, \ y, \ z$ の和である$20$個を〇で表し, \ 二つの仕切り｜を用いて$x, \ y, \ z$ の$3$つに分ける.

例えば$x=4, \ y=10, \ z=6$ のときは○○○○｜○○○○○○○○○○｜○○○○○○と表される.

このとき, \ $(x, \ y, \ z)$ の組は〇と｜の順列を考えればよく, \ $\dfrac{22!}{20! \ 2!} = \underline{231 \ (組)}$.

\bigskip
\bigskip

4.

(1) サイコロを$4$回投げたとき出る目の最大値が$4$となる確率を求めよ

\bigskip

最大値が$4$になるということは, \ サイコロを$4$回投げたときに少なくとも$1$回は$4$の目が出て, \ それ以外は$4$以下の目が出るということになる.

これはすべて$4$の目が出る確率からすべて$3$の目が出る確率を引いたものと考えられるため, \ 確率は

$\left(\dfrac{4}{6} \right)^4 - \left(\dfrac{3}{6} \right)^4
= \underline{\dfrac{175}{1{,}296}}$.

\bigskip
\bigskip

(2) 袋の中に白玉$4$個, \ 赤玉$5$個が入っており, \ AとBの2人がこの順に玉を取り出す(取り出した玉は元に戻さない).

Bが白玉を取り出したとき, \ Aの取り出した玉が赤であった確率を求めよ

\bigskip

Aが赤玉を取り出す事象を$X$, \ Bが白玉を取り出す事象を$Y$とする.

このとき, \ 求める確率は$P_Y X = \dfrac{P(X \cap Y)}{P(Y)}$.

$P(X \cap Y) = \dfrac{5}{9} \cdot \dfrac{4}{8} = \dfrac{20}{72}, \ \ 
P(\overline{X} \cap Y) = \dfrac{4}{9} \cdot \dfrac{3}{8} = \dfrac{12}{72}$ \ \ ($\overline{X}$ は$X$の余事象).

$P(Y) = P(X \cap Y) + P(\overline{X} \cap Y) = \dfrac{32}{72}$ であるから, \ $\underline{P_Y X = \dfrac{P(X \cap Y)}{P(Y)} = \dfrac{5}{8}}$.

\bigskip
\bigskip

(3) サイコロを$3$回投げ, \ 出た目の最大値を$X$とする. $X$の期待値$E(X)$ と分散$V(X)$ を求めよ

\bigskip

$X = 1, ..., 6$ となるときのそれぞれの確率($=P_X$ とする)を(1)と同様にして求める.

$X = 1$ のとき, \ $P_1 = \left(\dfrac{1}{6} \right)^3 = \dfrac{1}{6^3}$.

$X = 2$ のとき, \ $P_2 = \left(\dfrac{2}{6} \right)^3 - \left(\dfrac{1}{6} \right)^3 = \dfrac{7}{6^3}$.

$X = 3$ のとき, \ $P_3 = \left(\dfrac{3}{6} \right)^3 - \left(\dfrac{2}{6} \right)^3 = \dfrac{19}{6^3}$.

$X = 4$ のとき, \ $P_4 = \left(\dfrac{4}{6} \right)^3 - \left(\dfrac{3}{6} \right)^3 = \dfrac{37}{6^3}$.

$X = 5$ のとき, \ $P_5 = \left(\dfrac{5}{6} \right)^3 - \left(\dfrac{4}{6} \right)^3 = \dfrac{61}{6^3}$.

$X = 6$ のとき, \ $P_6 = \left(\dfrac{6}{6} \right)^3 - \left(\dfrac{5}{6} \right)^3 = \dfrac{91}{6^3}$.

よって, \ $E(X) = 1P_1 + 2P_2 + 3P_3 + 4P_4 + 5P_5 + 6P_6 = \underline{\dfrac{1{,}071}{6^3}}$.

また, \ $E(X^2) = \displaystyle \sum_{k=1}^6 k^2 P_k
= \dfrac{5{,}593}{6^3}$.

$V(X) = E(X^2) - (E(X))^2 = \dfrac{5{,}593}{6^3} - \left(\dfrac{1{,}071}{6^3} \right)^2 = \underline{\dfrac{61{,}047}{46{,}656}.}$

\end{document}
